\chapter{把黎曼写进 Lean：陈述、边界与诚实}
\label{chap:lean-riemann}

\begin{keybox}[本章要点]
mathlib v4.36.0 已经含有解析延拓后的 $\zeta$、函数方程、平凡零点，
甚至直接定义了 \texttt{RiemannHypothesis}。本章展示：
如何把百万美元难题写成一行 \texttt{def}，如何推导它的一个等价版本，
以及——更重要的——\emph{怎样诚实地"陈述"而不是假装"证明"}。
\end{keybox}

\section{mathlib 里已经有什么}
\label{sec:08-what-exists}

用 \texttt{\#check} 查一下（下面是真实输出）：

\begin{lstlisting}[language=Lean,numbers=none,frame=none,xleftmargin=0pt]
#check riemannZeta
-- riemannZeta (a : ℂ) : ℂ

#check completedRiemannZeta₀_one_sub
-- completedRiemannZeta₀_one_sub (s : ℂ) :
--     completedRiemannZeta₀ (1 - s) = completedRiemannZeta₀ s

#check zeta_eq_tsum_one_div_nat_cpow
-- zeta_eq_tsum_one_div_nat_cpow {s : ℂ} (hs : 1 < s.re) :
--     riemannZeta s = ∑' n : ℕ, 1 / (n : ℂ) ^ s

#check RiemannHypothesis
-- RiemannHypothesis : Prop
\end{lstlisting}

几点说明：

\begin{itemize}
  \item \texttt{riemannZeta : ℂ → ℂ} 是\emph{已经解析延拓到全平面}的 $\zeta$。
        它并不是用 $\sum n^{-s}$ 定义的（那个级数在 $\Ree s\le1$ 发散），
        而是用 Hurwitz zeta / Mellin 变换构造的；
        \texttt{zeta\_eq\_tsum\_one\_div\_nat\_cpow} 才说明在 $\Ree s>1$
        时它等于欧拉的级数。这正是 \cref{thm:02-continuation} 的形式化版本。
  \item \texttt{completedRiemannZeta₀} 类似第一编的 $\Lambda_0(s)$ 去掉极点后的整函数；
        上面的定理就是函数方程 $\Lambda_0(1-s)=\Lambda_0(s)$。
  \item \texttt{∑'} 是拓扑无穷和（\texttt{tsum}）；\texttt{(n : ℂ) \textasciicircum{} s}
        用了复数幂 \texttt{cpow}。这些在第一编里是"显然收敛"的东西，
        在类型里必须写清楚。
\end{itemize}

\section{函数方程：纸面与机器}
\label{sec:08-fe}

第一编 \cref{eq:02-fe3} 写成
\[
  \zeta(1-s)=2\,(2\pi)^{-s}\cos\!\left(\frac{\pi s}{2}\right)\Gamma(s)\,\zeta(s).
\]
mathlib 的 \texttt{riemannZeta\_one\_sub} 是（真实类型）：

\begin{lstlisting}[language=Lean,numbers=none,frame=none,xleftmargin=0pt]
theorem riemannZeta_one_sub {s : ℂ} (hs : ∀ n : ℕ, s ≠ -n) (hs' : s ≠ 1) :
    riemannZeta (1 - s) =
      2 * (2 * ↑Real.pi) ^ (-s) * Complex.Gamma s
        * Complex.cos (↑Real.pi * s / 2) * riemannZeta s
\end{lstlisting}

对照一下，会发现两件很有教育意义的事：

\begin{enumerate}[label=\arabic*.]
  \item \emph{附加条件}。纸面上常省略"避开 $\Gamma$ 的极点"这类条件；
        这里必须显式写成 \texttt{hs : ∀ n, s ≠ -n} 与 \texttt{hs' : s ≠ 1}。
        形式化迫使我们把定理的适用范围说完整。
  \item \emph{名称与记号的差异}。纸面写 $\pi$，机器写 \texttt{↑Real.pi}
        （把实数 $\pi$ 嵌入 $\C$）；纸面写 $\Gamma(s)$，机器写 \texttt{Complex.Gamma s}。
        这类"强制转换"在纸面上通常被省略，在类型里躲不掉。
\end{enumerate}

\section{把黎曼猜想写成一行 \texttt{def}}
\label{sec:08-rh}

mathlib 的原文（\texttt{Mathlib/NumberTheory/LSeries/RiemannZeta.lean}）是：

\begin{lstlisting}[language=Lean,numbers=none,frame=none,xleftmargin=0pt]
/-- A formal statement of the **Riemann hypothesis** - constructing a term
of this type is worth a million dollars. -/
@[wikidata Q205966]
def RiemannHypothesis : Prop :=
  ∀ (s : ℂ) (_ : riemannZeta s = 0) (_ : ¬∃ n : ℕ, s = -2 * (n + 1))
    (_ : s ≠ 1), s.re = 1 / 2
\end{lstlisting}

读法：对任意 $s$，若 $\zeta(s)=0$、且 $s$ 不是平凡零点
$-2,-4,-6,\dots$、且 $s\ne1$（极点），则 $\Ree s=1/2$。

\begin{histbox}[库作者的幽默]
那句 \texttt{constructing a term of this type is worth a million dollars}
不是玩笑话的夸张：克雷数学所确实为 RH 的证明悬赏一百万美元
\cite{clay2000}。而 \texttt{Lean} 里的"构造一个该类型的项"
恰恰就是"给出一个证明"。数学、逻辑与奖金在这里合成了一句话。
\end{histbox}

请注意 \texttt{RiemannHypothesis} 是一个 \texttt{Prop}，\emph{不是}带证明的数据：
\texttt{\#check RiemannHypothesis} 打印 \texttt{RiemannHypothesis : Prop}。
mathlib 里没有任何 \texttt{theorem RH\_proof : RiemannHypothesis}——
它只是被\emph{陈述}出来了。

\section{推导一个等价版本（真证明）}
\label{sec:08-critical}

教科书常用的说法是"临界带内的零点都在临界线上"。我们把这句话写成
\texttt{RH\_CriticalStrip}，并证明 \emph{mathlib 的 RH 蕴含它}：

\leanfile[caption={lean/Riemann/ZetaStatements.lean},label={lst:zetastat}]{../lean/Riemann/ZetaStatements.lean}

\texttt{RH\_implies\_criticalStrip} 的证明做了两件小事：

\begin{enumerate}[label=\arabic*.]
  \item 若 $s$ 是平凡零点 $s=-2(n+1)$，代入后
        \texttt{norm\_num} 把 $\Ree s$ 化简成 $2(n+1)<0$ 的形式，
        与前提 \texttt{hpos : 0 < s.re} 矛盾（\texttt{nlinarith} 收尾）。
  \item 若 $s=1$，与 \texttt{hlt : s.re < 1} 矛盾。
\end{enumerate}

\begin{warnbox}[为什么反过来难]
反方向（临界带版本 $\Rightarrow$ mathlib 版本）需要证明：
\emph{任何}非平凡零点都满足 $0<\Ree s<1$。这依赖第一编
\cref{thm:02-zeros} 中"$\Ree s=1$ 上无零点"这样的深层定理，
不是几行战术能完成的。所以文件里只证明了一个方向，并在注释里说清楚。
\emph{知道哪一步欠着，比假装证完更重要。}
\end{warnbox}

\section{如何"陈述"一个还没证明的东西}
\label{sec:08-honest}

形式化一个开放问题时，正确姿势是：\emph{精确陈述，留空证明}，
而且不偷偷引入 \texttt{axiom}。练习文件里给出一个例子：

\begin{lstlisting}[language=Lean,numbers=none,frame=none,xleftmargin=0pt]
def EulerProductHolds : Prop :=
  ∀ s : ℂ, 1 < Complex.re s →
    riemannZeta s = ∏' p : { p : ℕ // p.Prime },
                       1 / (1 - (p : ℂ) ^ (-s))
\end{lstlisting}

它把 \cref{thm:01-euler} 推广到复数 $s$ 并用无穷乘积 \texttt{∏'} 表达；
类型正确，但本项目没有给出证明（练习题第 4 题只要求\emph{陈述}）。
对比三种做法：

\begin{center}\small
\begin{tabular}{@{}p{0.34\textwidth}p{0.60\textwidth}@{}}
\toprule
\textbf{做法} & \textbf{评价}\\
\midrule
\texttt{def P : Prop := ...} & 正确：只定义命题\\
\texttt{axiom h : P} & 危险：把猜想当成公理，后续所有"证明"都不可信\\
\texttt{theorem t : P := by sorry} & 自欺：看起来是定理，其实有洞\\
\bottomrule
\end{tabular}
\end{center}

\begin{keybox}[形式化的伦理]
\texttt{sorry} 与 \texttt{axiom} 本身不是坏事——它们是显式声明的"欠条"。
坏事是把欠条藏起来，或者用 \texttt{axiom} 让一个假命题看起来被证明了。
\texttt{Lean} 的价值之一，就是让欠条\emph{无处可藏}。
\end{keybox}

\section{练习题}
\label{sec:08-exercises}

\leanfile[caption={lean/Riemann/Exercises.lean（故意含 sorry）},label={lst:exercises}]{../lean/Riemann/Exercises.lean}

六道题的难度与提示都写在注释里。做完后 \texttt{lake build} 不再打印
\texttt{sorry} 警告，就是通关标志。附录~\ref{app:exercises} 给出提示与部分答案。

\section{从这里往哪里走}
\label{sec:08-next}

\begin{description}[leftmargin=2.2em]
  \item[查库] \texttt{loogle.lean-lang.org}（按类型搜索定理）、
        \texttt{leansearch.net}、mathlib 文档\cite{mathlib4}、
        VS Code 的自动补全与 \texttt{\#check}。
  \item[读库] \texttt{Mathlib/NumberTheory/LSeries/RiemannZeta.lean}、
        \texttt{WeakPNT.lean}、\texttt{PrimeCounting.lean} 都值得读；
        你会发现形式化里的 $\zeta$ 用 Hurwitz zeta 与 Mellin 变换构造，
        与第一编的 $\theta$ 函数路线不同但等价。
  \item[教材] \emph{Mathematics in Lean}\cite{mil2023}、
        \emph{Theorem Proving in Lean 4}\cite{lean4ref}。
  \item[进阶] 试着把 \cref{thm:01-euler} 的有限证明推广成
        \texttt{EulerProductHolds} 的真证明（需要 Dirichlet 级数与收敛性的
        大量基础设施）；或者给 \texttt{Exercises.lean} 第 5 题补上
        "RH $\Rightarrow$ 误差项"的陈述。
\end{description}

\section{第二编总结}
\label{sec:08-summary}

\begin{center}\small
\begin{tabular}{@{}p{0.30\textwidth}p{0.64\textwidth}@{}}
\toprule
\textbf{你学到的} & \textbf{对应章节}\\
\midrule
命题即类型、可信内核、信任清单 & \cref{chap:lean-intro}\\
欧几里得定理的逐战术拆解 & \cref{sec:07-euclid}\\
有限欧拉乘积 = 唯一分解 + 计算 & \cref{sec:07-eulerproduct}\\
M\"obius 反演 = 卷积逆 & \cref{sec:07-moebius}\\
$\pi$、单调性、素数定理的\emph{陈述} & \cref{sec:07-pnt}\\
$\zeta$、函数方程、RH 的形式化 & \cref{chap:lean-riemann}\\
\bottomrule
\end{tabular}
\end{center}

最后一句：第一编里黎曼把素数变成零点，第二编里我们把零点变成类型。
这两件事其实共享同一个信念——\emph{把模糊的东西放到精确的语言里，
它才会真正被理解}。
